\documentclass[letterpaper]{article}
\usepackage[submission]{aaai2027}
\usepackage[hyphens]{url}
\usepackage{graphicx}
\urlstyle{rm}
\def\UrlFont{\rm}
\usepackage{natbib}
\usepackage{caption}
\usepackage{amsmath}
\usepackage{amssymb}
\usepackage{algorithm}
\usepackage{algorithmic}
\usepackage{booktabs}
\frenchspacing

\pdfinfo{
/TemplateVersion (2027.1)
}

\setcounter{secnumdepth}{0}

% note (luojiaxuan): Keep every unmeasured result explicit so the draft never
% silently turns an intended experiment into an empirical claim.
\newcommand{\pending}[1]{\textbf{[Pending: #1]}}
\newcommand{\method}{\textsc{CausalCache}}
\newcommand{\hf}{\mathrm{HF}}
\newcommand{\lf}{\mathrm{LF}}
\newcommand{\cost}{\operatorname{cost}}

\title{CausalCache: Budget-Conditioned Behavioral Restoration for Long-Horizon GUI Memory}
\author{Anonymous Submission}
\affiliations{}

\begin{document}

\maketitle

\begin{abstract}
Long-horizon GUI agents must decide which historical visual evidence to expose
to an action policy under a limited multimodal context budget. Existing memory
controllers commonly learn from recency, similarity, attention, or salience
labels, none of which directly measures whether high-fidelity evidence changes
the policy's executable decision. We introduce \method, a policy-grounded
supervision framework that upgrades one event at a time from a fixed-cost
summary to its archived high-fidelity representation and measures how much the
upgrade restores a validated frozen policy's action behavior. We formulate a
budget-conditioned, Shapley-style restoration value whose coalitions match the
deployment budget, then distill the expensive offline values into a
query-conditioned multi-budget gate. The method limits policy-visible context,
not persistent storage, and filters reference states whose full-history action
is not validated. We evaluate whether restoration supervision improves the
closed-loop success--memory frontier and whether retained restoration value
predicts success after controlling for next-action likelihood, context cost,
and task horizon. \pending{Insert the principal AndroidWorld result and the
matched-NLL effect size.}
\end{abstract}

\section{Introduction}

A GUI trajectory can contain screenshots, actions, text arguments, coordinates,
and state changes accumulated over dozens of decisions. Passing every event to
a multimodal policy is expensive and can amplify context rot, while replacing
all screenshots with summaries discards exact visual evidence. A memory
controller must therefore decide which historical events deserve high-fidelity
exposure at each decision.

Prior controllers select memory using recency, retrieval similarity, attention,
learned salience, or task-state heuristics. These signals need not agree with
the frozen action policy's actual dependence on an event. Two memories can even
have similar successor-action negative log-likelihood (NLL) while supporting
different closed-loop outcomes. This motivates a narrower question: \emph{how
much does upgrading a historical event from a deterministic low-fidelity record
to its high-fidelity visual evidence restore a validated policy's executable
action behavior?}

\method answers this question offline and uses the answers as supervision. It
maintains a raw event archive, a cheap summary/index, and a limited
policy-visible high-fidelity context. For validated decision states, it samples
near-budget mixed-fidelity coalitions and estimates each event's marginal
behavioral restoration. The estimator is explicitly Shapley-style; our claimed
contribution is the GUI intervention surface, the executable policy-behavior
value, its alignment with the deployment budget, and closed-loop validation,
not a new generic attribution estimator. A lightweight query-time gate then
predicts these values from the current state, event summary, optional cheap
visual embedding, age, and requested budget.

Our evaluation is organized around four falsifiable questions:
\begin{enumerate}
    \item Does a restoration oracle improve executable task success at a fixed
    policy-visible visual budget?
    \item Can a distilled gate approach the oracle without material online
    latency?
    \item Does restoration supervision outperform the same gate architecture
    trained on next-action NLL or shuffled labels?
    \item After matching or controlling for NLL, tokens, and horizon, does
    retained restoration value still explain terminal success?
\end{enumerate}

\section{Related Work}

\paragraph{GUI memory.}
MementoGUI learns multimodal write, compression, and retrieval operators for a
frozen GUI action model \cite{zeng2026mementogui}. AndroTMem retrieves anchored
intermediate states connected to the current subgoal \cite{shi2026androtmem}.
These systems establish the importance of learned controllers and
dependency-critical events. \method instead derives selector supervision from
the marginal behavior change induced by restoring the high-fidelity version of
the same event.

\paragraph{Interventional memory selection.}
CMI compares agent behavior under memory interventions
\cite{srivastava2026cmi}, while ATMem estimates how active memory states affect
GUI task completion \cite{liu2026atmem}. Our scope is intentionally narrower:
we hold the action policy fixed, intervene on low- versus high-fidelity GUI
evidence, assign credit near the actual context budget, and validate the dense
decision-time surrogate through executable closed-loop outcomes.

\paragraph{Shapley-style attribution.}
Permutation marginal contributions originate in cooperative game theory
\cite{shapley1953value}. We use this tool as an estimator, but change its value
function and coalition distribution to reflect mixed-fidelity GUI context and a
fixed operating budget.

\section{Problem Formulation}

Let the history before decision $t$ be
\begin{equation}
H_t = \{e_1,\ldots,e_{t-1}\}, \qquad
e_j=(o_{j-1},a_j,o_j,\delta_j),
\end{equation}
where $\delta_j$ is an optional deterministic UI delta. Every event has an
archived high-fidelity representation $e_j^{\hf}$ and a fixed-cost low-fidelity
record $\tilde e_j$. A query-time selector chooses the events exposed at high
fidelity:
\begin{equation}
M_t \subseteq H_t, \qquad
\sum_{j\in M_t} c_j \leq B,
\end{equation}
where $c_j$ is the visual-token cost and $B$ is the policy-visible context
budget. Unselected events remain available as low-fidelity records. Raw events
are retained in an external archive, so an event can be retrieved again at a
later query; $B$ is not a persistent storage or eviction budget.

The action policy $\pi_0$ is frozen. Only the selector is trained. Our goal is
to improve executable task success across budgets without attributing gains to
extra visual tokens or a stronger action policy.

\section{CausalCache}

\subsection{Two-Level Event Representation}

The deterministic low-fidelity record contains
\begin{quote}
\small\ttfamily
step\_id, action\_type, target\_text\_or\_coordinate\_bin,\newline
deterministic\_ui\_delta, result\_status.
\end{quote}
The archive stores the original screenshots, action, UI delta, and metadata.
The gate may additionally use a precomputed cheap visual embedding $z_j$; this
avoids asking the summary alone to predict the importance of information that
the summary necessarily omitted. The action policy receives mixed-fidelity
inputs through a normal forward pass; we do not splice cached keys and values
across incompatible contexts.

\subsection{Validated Full-History Reference}

For a full-history state, the frozen policy defines
\begin{equation}
q_t = \pi_0(\cdot\mid g,o_t,H_t).
\end{equation}
We generate attribution labels only when the full-history canonical action
$a_t^*$ is validated by an executor-compatible reference action or belongs to a
successful trajectory. We report reference coverage and compare against an
unfiltered teacher. Thus, full history is a reference condition, not an assumed
oracle.

\subsection{Executable Behavioral Distance}

For a high-fidelity coalition $S$, all events outside $S$ use their low-fidelity
records. We measure pathwise divergence along the canonical action rather than
claiming a full sequence-action KL:
\begin{equation}
D_t(S) = \alpha D_{\mathrm{type}}(S)
       + \beta D_{\mathrm{target}}(S)
       + \gamma D_{\mathrm{text}}(S).
\end{equation}
Action type uses teacher-forced token divergence; coordinates are mapped to UI
elements or coarse bins; equivalent swipe grammars are mapped to viewport
scroll direction; text arguments are normalized before comparison.
Executable action match is the primary action metric, with raw token KL as an
auxiliary analysis.

\subsection{Budget-Conditioned Restoration}

For event $j$, let $\mathcal{C}_B(j)$ be a distribution over feasible
coalitions $S\subseteq H_t\setminus\{j\}$ that leave room for $c_j$ and are
near budget. The budget-conditioned restoration value is
\begin{equation}
G_{j,t}^{(B)} =
\mathbb{E}_{S\sim\mathcal{C}_B(j)}
\left[D_t(S)-D_t(S\cup\{j\})\right].
\label{eq:restoration}
\end{equation}
When all high-fidelity blocks have equal cost, $\mathcal{C}_B(j)$ is uniform
over coalitions of size $B-1$. Unlike a standard Shapley value averaged over all
coalition sizes, Equation~\ref{eq:restoration} asks whether restoring $j$ helps
when the system already operates near its deployed budget. We use shared,
antithetic, or stratified coalition samples where applicable and report standard
errors, ranking stability, and top-$B$ overlap across seeds.

\begin{algorithm}[t]
\caption{Offline Budget-Conditioned Restoration Labels}
\label{alg:labels}
\begin{algorithmic}[1]
\REQUIRE Validated state $(g,o_t,H_t,a_t^*)$, budget $B$, samples $K$
\STATE Run full history once to cache reference behavior $q_t$
\FOR{$k=1$ to $K$}
    \STATE Sample a shared near-budget coalition design for all events
    \FOR{each event $j$ evaluated by the design}
        \STATE Run mixed-fidelity policy for $S$ and $S\cup\{j\}$
        \STATE Record $D_t(S)-D_t(S\cup\{j\})$
    \ENDFOR
\ENDFOR
\STATE Return means, standard errors, and sample metadata
\end{algorithmic}
\end{algorithm}

\subsection{Multi-Budget Query-Time Gate}

The distilled selector predicts
\begin{equation}
s_{j,t}=f_\theta(g,o_t,\tilde e_j,z_j,\Delta t,B)
\end{equation}
with regression, pairwise-ranking, and budget-aware set losses. At inference,
the selector solves a cost-constrained positive-value selection problem:
\begin{equation}
M_t=\arg\max_{M:\sum_{j\in M}c_j\le B}
\sum_{j\in M}\max(s_{j,t}-\tau,0).
\end{equation}
For equal-cost event blocks this reduces to selecting up to, rather than exactly,
the top $B$ events above threshold $\tau$. This permits a null decision when all
candidate restorations are predicted to be harmful or useless.

\section{Experimental Protocol}

\subsection{Benchmarks and Policies}

AndroidWorld \cite{rawles2025androidworld} is the sole primary closed-loop
benchmark. GUIOdyssey \cite{lu2025guiodyssey} or an equivalent long-horizon
mobile dataset provides offline attribution and action analysis. We use one
primary frozen policy and one transfer backbone. Every comparison shares the
same low-fidelity channel, action policy, gate architecture where applicable,
and visual-token accounting.

\subsection{Comparisons}

The compact main comparison includes summary only, recent, similarity, the same
gate trained on next-action NLL or salience, the strongest reproducible GUI
memory controller, a restoration oracle, and distilled \method. A random label
shuffle uses the same training states and gate capacity as a negative control.

\subsection{Success--Memory Frontier}

We report closed-loop task success against visual tokens and end-to-end latency,
plus success per unit memory/compute. Restoration must improve the Pareto
frontier rather than consume additional context. \pending{Insert main frontier
plot and paired confidence intervals.}

\subsection{Matched-NLL Mechanism Test}

Within task, budget, and horizon strata, we match memory decisions by successor
NLL and test whether retained restoration mass remains associated with terminal
success. The preregistered conditional model is
\begin{equation}
\begin{aligned}
\Pr(y=1)=\sigma\!\bigl(&\beta_0+\beta_1 R+\beta_2\mathrm{NLL}\\
&+\beta_3\mathrm{tokens}+\beta_4\mathrm{horizon}\bigr),
\end{aligned}
\end{equation}
where $R$ is retained restoration mass. The mechanism claim requires a stable
$\beta_1>0$ under paired bootstrap confidence intervals and reasonable matching
tolerances. We will report all eligible pairs rather than selected anecdotes.

\subsection{Attribution Cost and Stability}

We vary coalition samples $K$ and report attribution standard error,
Spearman rank correlation to a high-$K$ estimate, top-budget Jaccard overlap,
oracle success, total offline labeling compute, and online selector latency.
The per-permutation telescoping identity is used only as an implementation sanity
check. Coalition reconstruction error measures non-additive interactions.

\section{Limitations and Claim Boundary}

Restoration value is a dense decision-time surrogate. It is not a direct
long-horizon causal effect, and its relevance to long-horizon control must be
established through closed-loop success and the matched-NLL analysis. The causal
language in this paper refers only to a controlled low-/high-fidelity
intervention on a frozen policy's behavior. We do not identify environment-level
structural causality and do not learn a world model. The method should be
weakened or abandoned if restoration fails to beat recency, similarity, and
attention; if its matched-NLL association disappears; if the distilled gate
cannot approach the oracle; or if gains come from extra tokens.

\section{Conclusion}

\method turns mixed-fidelity behavioral restoration into policy-grounded
supervision for GUI memory selection. Its central empirical burden is not that
attribution can fit a selector, but that budget-aligned restoration identifies
high-fidelity evidence that improves executable long-horizon control beyond
what next-action likelihood and context size explain.

\bibliography{references}

\end{document}
